\chapter{Purpura}

\section*{Definition}
Purpura is visible, non-blanching extravasation of blood into the skin or mucous membranes, classified by size: petechiae (<2 mm), purpura (2--10 mm), and ecchymoses (>10 mm, bruises). It results from disorders of platelets, coagulation, or vessel walls.

\section*{What Happens in Your Body}
Purpura arises from either thrombocytopenia/platelet dysfunction (failure of primary hemostatic plug), coagulation factor deficiency (impaired fibrin clot), or vasculitic damage (inflammatory vessel wall disruption). Platelet-type purpura typically produces non-palpable petechiae on dependent areas; palpable purpura suggests leukocytoclastic vasculitis with immune complex deposition and neutrophil-mediated vessel injury.

\section*{Causes}
\begin{itemize}
  \item \textbf{Platelet disorders:} Immune thrombocytopenia (ITP), thrombotic thrombocytopenic purpura (TTP), drug-induced (heparin, quinine, sulfonamides), bone marrow failure, von Willebrand disease
  \item \textbf{Coagulopathies:} Hemophilia A/B, liver disease, vitamin K deficiency, disseminated intravascular coagulation (DIC), anticoagulant therapy (warfarin, DOACs)
  \item \textbf{Vascular:} Vasculitis (IgA vasculitis/Henoch--Sch\"onlein, ANCA-associated), senile purpura, corticosteroid-induced, Ehlers--Danlos syndrome, scurvy (vitamin C deficiency)
  \item \textbf{Infectious:} Meningococcemia, endocarditis, dengue, RMSF, leptospirosis, sepsis with DIC
\end{itemize}

\section*{When to See a Doctor}
See your doctor if:
\begin{itemize}
  \item Acute onset of widespread petechiae or purpura, especially in a febrile or ill-appearing child or adult
  \item Palpable purpura, especially on lower extremities (suggests vasculitis)
  \item Mucosal bleeding: epistaxis, gingival bleeding, hematuria, or hematochezia
  \item Rapidly spreading ecchymoses with hemodynamic instability or compartment syndrome
\end{itemize}

\section*{Related Symptoms}
\begin{itemize}
  \item Petechiae, ecchymoses, hematoma, epistaxis, gingival bleeding, hematuria, menorrhagia
\end{itemize}
\textbf{Important:} Non-palpable petechiae in a well-appearing afebrile person with isolated thrombocytopenia most likely represents ITP, whereas palpable purpura with fever, arthralgia, or abdominal pain warrants urgent vasculitis evaluation.

\section*{Self-Care / Home Management}
\begin{itemize}
  \item Avoid aspirin, NSAIDs, and other antiplatelet agents if purpura or bleeding tendency
  \item Apply ice and pressure to minor bleeding; avoid intramuscular injections
  \item Use a soft toothbrush and electric razor to minimize mucosal trauma
\end{itemize}

\section*{How Your Doctor Will Evaluate You}
\begin{itemize}
  \item CBC with platelet count, peripheral blood smear, PT/PTT/INR, fibrinogen, and D-dimer
  \item Bleeding time or platelet function assay if platelet count normal; von Willebrand panel if indicated
  \item Vasculitis workup: ESR, CRP, ANA, ANCA, cryoglobulins, complement levels, urinalysis, skin biopsy
\end{itemize}

\section*{Treatment Options}
\begin{itemize}
  \item ITP: Corticosteroids (prednisone 1 mg/kg/day), IVIG, thrombopoietin receptor agonists; splenectomy for refractory
  \item TTP: Urgent plasma exchange, high-dose corticosteroids, rituximab
  \item Coagulopathy: Factor replacement (hemophilia), vitamin K or FFP (warfarin reversal), stop offending anticoagulant
  \item Vasculitis: Treat underlying cause; corticosteroids and steroid-sparing immunosuppressants (cyclophosphamide, rituximab) for autoimmune vasculitis
  \item Discontinue offending medications in drug-induced purpura
\end{itemize}

\section*{References}
\begin{thebibliography}{99}
\bibitem{ref1} Neunert C, Lim W, Crowther M, et al. "The American Society of Hematology 2019 evidence-based practice guideline for immune thrombocytopenia." \textit{Blood Advances}, 2019;3(23):3829--3866.
\bibitem{ref2} Kwaan HC. "Thrombotic thrombocytopenic purpura and other thrombotic microangiopathies." \textit{New England Journal of Medicine}, 2017;376(12):1167--1176.
\bibitem{ref3} Jennette JC, Falk RJ, Bacon PA, et al. "2012 revised International Chapel Hill Consensus Conference Nomenclature of Vasculitides." \textit{Arthritis \& Rheumatism}, 2013;65(1):1--11.
\bibitem{ref4} George JN, Nester CM. "Syndromes of thrombotic microangiopathy." \textit{New England Journal of Medicine}, 2014;371(7):654--666.
\bibitem{ref5} O'Brien SH, Smith KJ, Moffett BS, et al. "Purpura: Evaluation and management in children and adults." \textit{The Lancet Haematology}, 2020;7(4):e342--e351.
\bibitem{ref6} Piaggio G, Elbourne DR, Pocock SJ, et al. "Reporting of noninferiority and equivalence randomized trials: An extension of the CONSORT statement." \textit{JAMA}, 2006;295(10):1152--1160. [Vasculitis biopsy criteria]
\end{thebibliography}
